% $Author: stef $
% $Date: 2008-04-04 17:14:31 +0200 (Fri, 04 Apr 2008) $
% $Revision: 318 $
%=================================================================
\ifx\wholebook\relax\else
% --------------------------------------------
% Lulu:
    \documentclass[a4paper,10pt,twoside]{book}
    \usepackage[
        papersize={6in,9in},
        hmargin={.75in,.75in},
        vmargin={.75in,1in},
        ignoreheadfoot
    ]{geometry}
    \input{../common.tex}
    \pagestyle{headings}
    \setboolean{lulu}{true}
% --------------------------------------------
% A4:
%   \documentclass[a4paper,11pt,twoside]{book}
%   \input{../common.tex}
%   \usepackage{a4wide}
% --------------------------------------------
    \graphicspath{{figures/} {../figures/}}
    \begin{document}
%   \renewcommand{\nnbb}[2]{} % Disable editorial comments
    \sloppy
\fi

\chapter{}\label{cha:}

%%%%%%%%%%%%%%%%%%%%%%%%%%%%%%%%%%%%%%%%%%%%%%%%%%%%%%%%%%%%%%%%%%%%%%%%
\begin{figure}[h]
	\centerline{\includegraphics{}}
	\caption{Ready-to-use files for the Macintosh.  
	\label{fig:macfiles}}
\end{figure}

%%%%%%%%%%%%%%%%%%%%%%%%%%%%%%%%%%%%%%%%%%%%%%%%%%%%%%%%%%%%%%%%%%%%%%%%
\begin{exofigwithsize}[0.5]{\includegraphics[width=3cm]{turtleMSmallStairs}}{A Staircase}\label{xp:letterA}
You are not limited in your robot drawings to squares. You can create a wide range of geometrical figures.
For example, here is a drawing of a small staircase. Write a script to reproduce this drawing. 
\end{exofigwithsize}

\begin{exonofigtitle}{Moving Clock Hands}
Experiment with different angle values for each of the two robots; that is,change the angle values for the two turn 
methods. Then, compare the effect of the method \ct{turnLeft: 60} (for pica) and \ct{turnRight: 300} (for daly). 
You can see that turning left 60 degrees yields the same result as turning right 300 degrees. This is so because 
the sum of the two values is 360 degrees, that is, a full circle. 
\end{exonofigtitle}

%%%%%%%%%%%%%%%%%%%%%%%%%%%%%%%%%%%%%%%%%%%%%%%%%%%%%%%%%%%%%%%%%%%%%%%%
\begin{scriptfigwithsize}[0.4]{\includegraphics[width=5cm]{turtleMLetterA}}{The letter A}\label{scr:letterA}
	| pica | 
	pica := Bot new. 
	pica north. 
	pica go: 100. 
	pica east. 
	pica go: 100. 
	pica south. 
	pica go: 100. 
	pica north. 
	pica go: 50. 
	pica west. 
	pica go: 100
\end{scriptfigwithsize}

\begin{script}[myster1]{What does pica do? (Problem 1)}
	| pica | 
	pica := Bot new. 
	pica go: 100. 
	pica turnLeft: 45. 
	pica go: 50. 
	pica turnLeft: 45. 
	pica go: 100 
\end{script}



\ifx\wholebook\relax\else
    \end{document}
\fi

%%% Local Variables:
%%% coding: utf-8
%%% mode: latex
%%% TeX-master: t
%%% TeX-PDF-mode: t
%%% ispell-local-dictionary: "english"
%%% End:
